\documentclass[a4paper, 12pt]{article}
\usepackage[utf8]{inputenc}
\usepackage[T2A]{fontenc}
\usepackage[russian]{babel}
\usepackage{geometry}
\geometry{left=20mm,right=20mm,top=20mm,bottom=20mm}
\usepackage{enumitem}
\usepackage{xcolor}
\usepackage{amsmath}

\title{Возможные вопросы комиссии и ответы}
\author{Vehicle Re-ID: ResNet-50 / ResNet-101 vs ViT-Small}
\date{}

\begin{document}
\maketitle

\section*{Вопросы и ответы}
% ============================================

\subsection*{Теоретические вопросы}

\textbf{В1: Почему ViT показывает лучшие результаты, чем ResNet?}

\textit{Ответ:} ViT использует механизм self-attention, который с первого слоя имеет глобальное рецептивное поле. Это позволяет модели:
\begin{itemize}
  \item Учитывать зависимости между удалёнными частями изображения (например, связать форму фар и заднего бампера)
  \item Лучше моделировать контекст всего автомобиля
  \item Эффективнее использовать pretrain на ImageNet-21k (14M изображений vs 1.2M у ImageNet-1k)
\end{itemize}
CNN строит иерархию признаков от локальных к глобальным, и глобальный контекст появляется только в последних слоях. Это подтверждается количественно: коэффициент силуэта ViT-Small (0.260) в 10 раз превышает аналогичный показатель ResNet-50 (0.024), а межкластерное расстояние у ViT-Small в 1.54 раза больше.

\vspace{1em}
\textbf{В2: Что такое BNNeck и зачем он нужен?}

\textit{Ответ:} BNNeck (Batch Normalization Neck) --- архитектурный приём, разделяющий пространства признаков для разных функций потерь:
\begin{itemize}
  \item \textbf{До BN}: признаки для Triplet Loss (метрическое обучение)
  \item \textbf{После BN}: признаки для Cross-Entropy (классификация)
\end{itemize}
Это уменьшает конфликт целей между потерями и стабилизирует обучение. На инференсе используется L2-нормализованное представление после BN.

\vspace{1em}
\textbf{В3: Как работает Triplet Loss?}

\textit{Ответ:} Triplet Loss обучает модель так, чтобы расстояние между якорем (anchor) и позитивом (тот же ID) было меньше, чем между якорем и негативом (другой ID), с зазором margin:
$$\mathcal{L} = \max(0, d(a,p) - d(a,n) + m)$$
где $d$ --- евклидово расстояние, $m=0.3$ --- margin. Hard mining выбирает самые сложные триплеты в батче для эффективного обучения.

\vspace{1em}
\textbf{В4: Что такое PK-семплинг?}

\textit{Ответ:} PK-семплинг --- стратегия формирования батча: выбираем P случайных идентичностей, для каждой берём K изображений. В работе ResNet-50: P=16, K=4 (batch=64); ViT-Small: P=12, K=4 (batch=48).

Преимущества:
\begin{itemize}
  \item Гарантируем наличие позитивных пар в батче
  \item Можем применить hard mining внутри батча
  \item Баланс между разнообразием ID и примерами каждого
\end{itemize}

\vspace{1em}
\textbf{В5: Что такое mAP и чем отличается от Rank-1?}

\textit{Ответ:}
\begin{itemize}
  \item \textbf{Rank-1 (CMC@1)}: доля запросов, где правильный ответ на первом месте
  \item \textbf{mAP}: средняя точность по всем релевантным результатам
\end{itemize}
Rank-1 = 90\% означает, что в 90\% случаев первый результат правильный.

mAP учитывает все правильные результаты и их позиции. Если одно транспортное средство присутствует в галерее 5 раз (с разных камер), mAP показывает, насколько хорошо все 5 ранжируются. ViT-Small: Rank-1 = 90.05\%, mAP = 69.41\%.

\subsection*{Практические вопросы}

\textbf{В6: Почему не использовали номерные знаки?}

\textit{Ответ:} По нескольким причинам:
\begin{itemize}
  \item Номера часто не видны: ракурс сбоку, блики, грязь
  \item Приватность: в некоторых сценариях использование номеров ограничено
  \item Научный интерес: показать, что ReID возможен по внешнему виду
  \item Практическое применение: камеры с низким разрешением
\end{itemize}

\vspace{1em}
\textbf{В7: Почему ResNet-101 показал результаты хуже, чем ResNet-50?}

\textit{Ответ:} Результат неочевидный: ResNet-101 (43.7M параметров) уступил ResNet-50 (24.7M): Rank-1 86.7\% против 88.7\%, mAP 57.7\% против 64.3\%. Две причины:
\begin{enumerate}
  \item VeRi-776 (37\,746 изображений, 575 ID) недостаточен для эффективной оптимизации 43.7M параметров --- перефиттинг верхних слоёв.
  \item При одинаковых гиперпараметрах (lr, epochs) более глубокая модель получает меньше шагов оптимизации «на параметр» --- недообучение нижних слоёв.
\end{enumerate}
Косвенно это подтверждает бо\'{о}льший прирост от re-ranking у ResNet-101 (+2.4\% mAP) --- пространство эмбеддингов менее структурировано.

\vspace{1em}
\textbf{В8: Как работает re-ranking и почему ViT выигрывает от него больше?}

\textit{Ответ:} K-reciprocal re-ranking учитывает взаимность ближайших соседей: если A входит в топ-$k_1$ соседей B, а B --- в топ-$k_1$ соседей A, они считаются взаимными соседями и получают дополнительный бонус близости. Параметры: $k_1=20$, $k_2=6$, $\lambda=0.3$.

Прирост mAP: ResNet-50 +3.2\%, ResNet-101 +2.4\%, ViT-Small +5.1\%. ViT выигрывает больше, потому что его эмбеддинги образуют более компактные и чёткие кластеры --- взаимность ближайших соседей работает лучше, когда соседи действительно «правильные».

\vspace{1em}
\textbf{В9: Что нашёл анализ по камерам (camera heatmap)?}

\textit{Ответ:} Для каждой пары (камера запроса, камера галереи) вычислялся Rank-1. Ключевые находки:
\begin{itemize}
  \item Лучшая пара C11$\to$C06: Rank-1 = 100\% (ResNet-50) --- камеры рядом, условия схожи.
  \item Проблемная камера C18: Rank-1 = 0--2.3\% для всех пар --- кардинально иные условия съёмки.
  \item ViT-Small превосходит ResNet-50 на 94.2\% ненулевых пар.
\end{itemize}
В реальной системе камера C18 --- сигнал к техническому обслуживанию. Суммарные метрики (Rank-1, mAP) это не выявляют.

\vspace{1em}
\textbf{В10: Зачем нужен YOLOv11 в системе ReID?}

\textit{Ответ:} VeRi-776 содержит уже вырезанные изображения транспортных средств (GT-кропы). Но в реальной жизни камера видит весь перекрёсток. YOLOv11n детектирует все ТС на кадре, вырезает bbox, после чего ReID-модель работает с отдельным кропом.

Влияние детектора на качество: Rank-1 падает с 90.05\% (GT-кроп) до 85.88\% (YOLO-кроп) --- это «цена реализма». 21.5\% запросов YOLOv11 не детектировал ТС (fallback на полное изображение). Для уменьшения этого разрыва можно дообучить YOLO на VeRi.

\vspace{1em}
\textbf{В11: Зачем расстояние Махаланобиса дало деградацию –31\% mAP для ViT?}

\textit{Ответ:} ReID-эмбеддинги после L2-нормализации лежат на единичной гиперсфере $\mathcal{S}^{383}$. На гиперсфере косинусное расстояние (= L2 для нормированных векторов) геометрически корректно. Матрица Махаланобиса $\Sigma^{-1}$ описывает эллипсоид в евклидовом пространстве --- «растягивание» пространства разрушает угловые отношения точек на сфере. Результат: ранжирование деградирует. Для ResNet-50 использовалось диагональное приближение --- эквивалентно нормировке по стандартному отклонению признаков, эффект практически нулевой (+0.28\% mAP).

\vspace{1em}
\textbf{В12: Что показывают карты внимания ViT-Small?}

\textit{Ответ:} Attention maps агрегируются из последних 4 трансформерных блоков (максимум по головам). При верных совпадениях Rank-1 модель концентрирует внимание на дискриминативных деталях: номерной знак, рисунок диска, спойлер, уникальные элементы кузова. При ошибках --- внимание размыто по форме и цвету. Важно: модель не обучалась с разметкой частей автомобиля --- она сама научилась смотреть на нужные места.

\vspace{1em}
\textbf{В13: Как масштабировать систему на миллионы изображений?}

\textit{Ответ:} Использовать приближённый поиск ближайших соседей:
\begin{itemize}
  \item \textbf{FAISS}: библиотека от Facebook для быстрого поиска векторов
  \item \textbf{IVF + PQ}: кластеризация + квантование векторов
  \item \textbf{HNSW}: графовый индекс для быстрого поиска
\end{itemize}
Для real-time поиска --- FAISS с косинусным расстоянием (миллисекунды). Re-ranking ($O(N^2\log N)$) --- только для offline-задач (найти машину из архива за день).

\subsection*{Вопросы на понимание кода}

\textbf{В14: Почему в ResNet изменили stride в layer4?}

\textit{Ответ:} В коде:
\begin{verbatim}
resnet.layer4[0].downsample[0].stride = (1, 1)
resnet.layer4[0].conv2.stride = (1, 1)
\end{verbatim}
Это убирает последний downsampling, увеличивая пространственное разрешение feature map с $8\times4$ до $16\times8$. Для ReID важны мелкие детали (номер, диски), и большее разрешение помогает их сохранить.

\vspace{1em}
\textbf{В15: Зачем используется GradScaler?}

\textit{Ответ:} GradScaler --- часть mixed precision training (FP16):
\begin{itemize}
  \item Ускоряет обучение в 2--3 раза
  \item Уменьшает потребление памяти GPU (критично при 6 GB VRAM GTX 1660 Super)
  \item Scaler масштабирует градиенты, предотвращая underflow в FP16
\end{itemize}

\subsection*{Каверзные вопросы}

\textbf{В16: Ваши результаты хуже SOTA. Почему это исследование ценно?}

\textit{Ответ:} Ценность работы:
\begin{enumerate}
  \item \textbf{Воспроизводимость}: полный код, понятный пайплайн, открытый датасет
  \item \textbf{Сравнение архитектур}: CNN vs Transformer на одинаковых условиях --- честный бенчмарк
  \item \textbf{Негативные результаты}: Mahalanobis деградирует на нормированных эмбеддингах; ResNet-101 хуже ResNet-50 --- это ценные практические выводы
  \item \textbf{End-to-end анализ}: от детекции YOLOv11 до визуализации attention maps и анализа по камерам
\end{enumerate}
SOTA методы используют ViT-Base (86M параметров), ensemble, дополнительные датасеты. Мы показали воспроизводимый baseline.

\vspace{1em}
\textbf{В17: А если автомобили одинакового цвета и модели?}

\textit{Ответ:} Это действительно сложный случай. Attention maps показывают: при ошибках модель фокусируется на цвете и форме, упуская уникальные детали. Модель опирается на:
\begin{itemize}
  \item Мелкие различия: царапины, наклейки, рисунок дисков
  \item Уникальные детали (если видны): государственный номер, логотип
  \item Особенности освещения и отражений
\end{itemize}
В таких случаях система выдаёт несколько кандидатов для ручной проверки. Анализ хитмап по парам камер показал конкретные проблемные комбинации (например, камера C18).

\vspace{1em}
\textbf{В18: Почему не использовали более новые методы: CLIP, Swin, DINOv2?}

\textit{Ответ:} Сознательный выбор baseline: показать сравнение CNN и Transformer на одинаковых условиях. Замена ViT-Small на ViT-Base с ImageNet-21k (TransReID, 2021) исторически сокращает разрыв с SOTA до 5--7\%, CLIP-претрейн закрывает его за счёт семантического выравнивания. Это упомянуто в разделе «Направления дальнейших исследований».

\end{document}
